% =====================================================================
% Conclusion , final, 2026-08-13.
%
% Every clause corresponds to a measurement. The submitted version concluded
% that encoding was "the dominant determinant" and that LeakyParallel/latency
% offered the strongest operating point; the strict evaluation supports a
% narrower and better-evidenced statement, and the difference is the paper.
% =====================================================================

\section{Conclusion and Future Work}

We presented a controlled design-space study of spiking neural networks for
network intrusion detection, conducted entirely within \texttt{snnTorch}:
\FNScrVariants{} configurations of nine neuron models and three spike encodings,
evaluated across four benchmark datasets over five seeds, with ranking rules
fixed in advance and a leakage-resistant protocol used to confirm the selection.

The complete screening identifies \texttt{LeakyParallel} with latency encoding
as the top-ranked configuration, while the stricter confirmation independently
places latency first on the encoding axis (mean rank \FNEncRank{} over
\FNEncBlocks{} blocks) and \texttt{LeakyParallel} first on the neuron axis
(\FNNeuRank{} of \FNNeuLevels{} families) by aggregate mean rank. Neither lead
is uniform. \FNScrWithinCD{} of \FNScrVariants{} screening configurations lie
within one critical difference of the leader; on the neuron axis the margin is
narrow enough that another family wins more individual blocks; and on the most
difficult scenario-disjoint CTU-13 fold both the preferred encoding and the
preferred neuron family change. Difficult scenario shifts can therefore reverse
design preferences that hold in aggregate.

Latency coding's practical advantage is not that it reaches decisions earlier.
It attains its terminal operating point sooner than rate coding on two of five
strict protocols and later on three, so temporal convergence is
dataset-dependent. What generalises is compute demand: rate encoding requires
$\FNSopRatioMin\times$ to $\FNSopRatioMax\times$ more synaptic operations per
sample on every strict protocol. The decomposition explains why a raw spike
count is an incomplete efficiency measure: rate emits only
$\FNTotRatioMin\times$ to $\FNTotRatioMax\times$ more spikes in total, but its
excess activity sits in the input layer, ahead of the high-fanout projection
where each spike is most expensive.

Three further analyses bound the result. Removing exact train--test duplicates,
which affect $\FNDupKdd\%$ of the official KDDCup99 test set, leaves the
encoding-axis leader and complete ordering intact, so the ordering is not a
consequence of memorisation. Detection falls from $\FNKnownLo$--$\FNKnownHi$ on
known attack types to $\FNUnseenLo$--$\FNUnseenHi$ on types absent from
training, for every evaluated configuration, so none of them addresses open-set
detection. And an operating threshold calibrated to $1\%$ false alarms on
validation data transfers well on CIC-IDS2017 but reaches $\FNFarTestNsl\%$ on
NSL-KDD and $\FNFarTestCtuB\%$ on the hardest CTU-13 fold, so calibrated
operating points cannot be assumed stable under distribution shift.

Taken together, these results suggest that SNN configuration \emph{rankings} are
considerably more stable than the absolute behaviour of the resulting detectors.
The recommended operating point survives an aggressive audit (train-only
preprocessing, seed-independent partitions, group-respecting splits, and
duplicate and class-support auditing), while detection rates, false-alarm
calibration and even encoding preference remain sensitive to the evaluation
distribution.

Three directions follow. First, the scenario-level inversion we observe on
CTU-13 invites a mechanistic study of when temporal-contrast coding becomes
preferable to timing-based coding, which our static per-flow formulation cannot
answer. Second, threshold-transfer instability under distribution shift is a
deployment problem in its own right and motivates calibration methods that do
not assume a stable benign distribution. Third, a like-for-like comparison
against conventional non-spiking detectors under identical preprocessing remains
open; it is a different question from the one studied here, and answering it
credibly requires its own fairness protocol.
